\documentclass[11pt]{article}
\usepackage[margin=0.65in]{geometry}
\usepackage{lmodern}
\usepackage{microtype}
\usepackage{enumitem}
\usepackage{tabularx}
\usepackage[table]{xcolor}
\usepackage{titlesec}
\usepackage[hidelinks]{hyperref}
\setlength{\parindent}{0pt}
\setlength{\parskip}{0.28em}
\setlist[itemize]{leftmargin=1.3em,itemsep=0.08em,topsep=0.1em}
\titleformat{\section}{\large\bfseries}{}{0pt}{}
\titlespacing*{\section}{0pt}{0.65em}{0.2em}
\pagestyle{plain}

\newcommand{\blankline}{\par\vspace{0.1em}\noindent\rule{\linewidth}{0.4pt}\par\vspace{0.35em}}
\newcommand{\smallblank}{\rule{0.43\linewidth}{0.4pt}}

\begin{document}

\begin{center}
{\Large\bfseries U1L6SP --- Short Paper 1}\par
{\LARGE\bfseries Prophecy Is Not Proof?}\par
{\large Week 6: Validity, Truth, and Soundness}
\end{center}

\noindent\textbf{Name:} \smallblank\hfill
\textbf{Date:} \rule{1.35in}{0.4pt}

\begingroup\small
\vspace{0.25em}
\fcolorbox{black}{gray!8}{%
\begin{minipage}{0.94\linewidth}
\textbf{The question}

In \textit{Macbeth}, the witches tell Macbeth that he ``shalt be king hereafter.'' Does this prophecy give Macbeth a sound reason to take action toward becoming king, or does it only tell him what may happen? Make and defend your own judgment.

You may answer \emph{yes}, \emph{no}, or \emph{partly}. A successful paper is not the one that guesses the teacher's answer. It is the one whose conclusion is earned by its reasons and evidence.
\end{minipage}}

\section*{Text support}

\begin{quote}
\textbf{The witches:} ``All hail, Macbeth, that shalt be king hereafter!'' (1.3)

\textbf{Macbeth:} ``If chance will have me king, why, chance may crown me, / Without my stir.'' (1.3)
\end{quote}

The accompanying \textit{U1L6SP Text Supplement} prints longer selections from Acts 1.3 and 1.7. These are lines worth focusing on, not limits on your evidence. You may use any accurate line, action, choice, image, or consequence from the play; you do not need outside research.

\section*{The assignment}

Write a short paper of \textbf{one to two handwritten pages} on separate lined paper. Give it the title \textit{Prophecy Is Not Proof?} Your paper should:

\begin{itemize}
\item answer the question with a clear claim;
\item explain the reasoning that could connect the prophecy to action;
\item test whether the conclusion really follows, even if the prophecy is true;
\item use at least two brief quotations or exact details from the supplement or any other part of \textit{Macbeth};
\item end with a careful judgment, including any limit or qualification your answer needs.
\end{itemize}

You do not need a five-paragraph essay. A clear beginning, a developed middle, and an honest ending are enough.

\section*{The closed-world rule}

This paper must come from \textbf{you}. You may use only your copy of \textit{Macbeth}, course handouts, and notes you made during class.

\textbf{Do not use AI, the internet, online summaries, outside criticism, a tutor, a parent, a friend, or anyone else's sentences or ideas.} Do not ask another person to plan, revise, or correct the paper. If you are uncertain, write through the uncertainty yourself. That struggle is part of the assignment.

At the end of your paper, write and sign: \emph{I wrote this paper without outside help.}

\section*{What counts on this first paper --- 10 points}

\rowcolors{2}{gray!8}{white}
\begin{tabularx}{\linewidth}{>{\bfseries}p{0.31\linewidth}Xr}
Claim & Gives a definite answer to the question. & 2\\
Logical work & Identifies and tests the reasoning between prophecy and action. & 3\\
Textual evidence & Uses at least two accurate quotations or details. & 2\\
Explanation & Connects the evidence to the claim instead of merely retelling the plot. & 2\\
Completion & Is legible, handwritten, and within the assigned length. & 1\\
\end{tabularx}

\vfill
\textit{This is a practice paper. Honest reasoning matters more than polish. Spelling and grammar will not lower the score unless they prevent the reader from understanding the argument.}
\endgroup

\newpage

\begin{center}
{\Large\bfseries Planning Scaffold}\par
{\large U1L6SP --- Use this page before writing}
\end{center}

\textbf{This page helps you think; it does not choose your answer.} Notes and fragments are enough. Your final paper goes on separate lined paper.

\section*{1. Your present answer}

In one sentence, what do you currently think the prophecy gives Macbeth a reason to do---or not to do?
\blankline
\blankline

\section*{2. Reconstruct the reasoning}

What claim does the prophecy make?
\blankline

What action or conclusion might Macbeth draw from it?
\blankline

What unstated bridge would be needed to move from the prophecy to that action?
\blankline
\blankline

\section*{3. Run both Week 6 tests}

\textbf{Form test:} If the relevant claims were true, would the proposed action or conclusion have to follow? Why or why not?
\blankline
\blankline

\textbf{Fact test:} Which claim, interpretation, or assumption most needs proof?
\blankline
\blankline

\section*{4. Choose your evidence}

\textbf{Quotation or exact detail 1:}
\blankline
\textbf{What it helps you show:}
\blankline

\textbf{Quotation or exact detail 2:}
\blankline
\textbf{What it helps you show:}
\blankline

\newpage

\section*{5. Make the judgment careful}

What might a reasonable reader say against your answer?
\blankline

What word such as \emph{because}, \emph{unless}, \emph{although}, \emph{only}, or \emph{partly} could make your claim more accurate?
\blankline

\section*{6. Simple writing path}

\begin{enumerate}[leftmargin=1.55em,itemsep=0.15em]
\item State your answer and name the reasoning problem.
\item Explain the possible bridge from prophecy to action.
\item Test that bridge with the play's evidence.
\item Give your qualified final judgment.
\end{enumerate}

\section*{7. Try your opening}

Draft two or three possible opening sentences. State your present answer and name the reasoning you will test. You may change these sentences when you write the paper.
\blankline
\blankline
\blankline

\section*{8. Details not to forget}

Record any final words, quotations, or ideas you want beside you while writing.
\blankline
\blankline
\blankline

\textbf{Before submitting:} Underline your main claim once. Circle the two places where you use evidence from the play.

\end{document}
